\clearpage
\section{Observational context congruence at equal input depths}
\label{sec:copying-term-congruence}
Unrestricted addressed filling now yields an observational congruence theorem
for every original annotated term context, provided the two input terms have
equal minimal translation depths. The conclusion is the existing
all-allowances observational equivalence of the original filled terms. It is
not merely equivalence of two graph fillings. The execution must respect
recursive graph isomorphism; no execution padding-invariance hypothesis is
needed for this equal-depth result.

The restriction is material. Lean checks that the unchanged private-cell
program of Proposition~12.2 has minimal translation depth three, whereas the
constant-result program has depth two. This section therefore does not settle
congruence for that pair, and does not prove their equivalence. The precise
remaining structural obligation is described below.

\subsection{Congruence through padded open graph contexts}

An open graph context has local allowance $r$ and host allowance $D$.
Padding its graph and selected address by $m$ steps yields allowances $r+m$
and $D+m$. Sole-hole and binding-free preservation supply the corresponding
certificates, so the result remains an admissible graph context. The iterated
filling law gives a recursive isomorphism between padding an original filling
and filling this padded context with the padded argument. Neither the outer
interface nor the selected local interface changes.


\lean{Copying/PaddedContextCongruence}{8}{15}{PaddedContextCongruence.pad}

\lean{Copying/PaddedContextCongruence}{16}{23}{PaddedContextCongruence.pad_fill_iso}

For any open program, padding first by $n$ and then by $m$ equals padding it
directly to the final allowance $(d+n)+m$. Rewrite both incremental paddings
as padding to a specified allowance, apply the exact graph-level transitivity
law, and use equality of the underlying graphs to identify the packaged
programs. This law includes all admissibility certificates by proof irrelevance.


\lean{Copying/PaddedContextCongruence}{24}{35}{PaddedContextCongruence.pad_padTo}

Suppose two fragments are observationally equivalent at every padded local
allowance. Their fillings in any fixed certified graph context are then
equivalent at every padded host allowance. Fix an increment $m$. The padded
filling isomorphisms give observational equivalences between each padded
original filling and its filling in the padded context, using the execution's
isomorphism law. The hypothesis supplies equivalence of the padded input
fragments. Apply graph-context congruence in the padded context, then compose
with the two surrounding equivalences, reversing the second one.

This argument does not claim that padding preserves execution of an individual
program. Instead, it uses an input comparison already supplied at the required
padded allowance. Thus the execution's isomorphism obligation suffices here.


\lean{Copying/PaddedContextCongruence}{36}{48}{PaddedContextCongruence.equiv_context_padded}

A scoped isomorphism also yields recursive graph isomorphism after padding
both graphs to any common adequate allowance. Unpack the two graph depths and
the scoped depth equality; eliminate that equality. The two graphs now have
the same index, so the previously proved common-allowance isomorphism congruence
applies. This transports structural filling correspondence to the allowance
at which an observation is requested.


\lean{Copying/PaddedContextCongruence}{49}{58}{PaddedContextCongruence.padTo}

\subsection{From actual filling to term-context congruence}
\termdef{copying-equal-depth-congruence}

Padding a graph to its existing allowance is the identity. This is the
zero-increment instance of the exact padding-to-addition law. Consequently,
interpreting a term context with the translated-program wrapper at the term's
own minimal depth gives exactly the original translation of the filled term.
Eliminate the stated equality of the chosen depth and the term's depth, remove
the identity padding, and apply the original constructor-interpretation law.
No term or distributor is simplified.


\lean{Copying/TermContextCongruence}{7}{13}{TermContextCongruence.padTo_self}

\lean{Copying/TermContextCongruence}{14}{25}{TermContextCongruence.apply_program_at_depth}

Let $t$ and $u$ have the same minimal translation depth $d$, and suppose
$t$ and $u$ are observationally equivalent in the existing all-allowances
sense. For every original context $C$, the theorem proves
\[
 \mathsf{ObservationalEquiv}(C[t],C[u]),
\]
provided the execution respects recursive isomorphism. All finite variable
and atom environments are retained; the hole may use variables bound by
surrounding lambdas.

\paragraph{One shared addressed context.}
Apply unrestricted graph filling to $C$ at allowance $d$. It supplies one
certified graph context $D$, with host allowance $q$ and local allowance
$d+n$, whose filling law works for every argument graph at allowance $d$.
Take $G$ and $H$ to be the original translated-program wrappers for $t$ and
$u$ at $d$. Their equality of minimal depths permits both exact
constructor-interpretation identities. The universal filling law therefore
relates $D[G\text{ padded by }n]$ to the original translation of $C[t]$, and
similarly for $u$. Each relation includes equality of the filled term's depth
with $q$.

\paragraph{Input comparison at every needed allowance.}
For an additional increment $m$, the two input programs are $G$ and $H$
padded first by $n$ and then by $m$. The exact program-padding law and the
translated-program coherence law identify these with the original translations
of $t$ and $u$ at allowance $(d+n)+m$. The assumed term equivalence supplies
precisely this comparison. No preservation of individual executions through
padding is used.

\paragraph{Every output observation.}
Fix any adequate output allowance $N$ for the two filled terms. The scoped
filling laws imply $q\leq N$, so write $N=q+m$. The padded graph-context
congruence theorem gives equivalence of the two filled graphs padded by $m$.
Pad each structural filling isomorphism to $N$. The resulting isomorphisms
identify those two padded fillings with the original translated-program
wrappers for $C[t]$ and $C[u]$ at $N$. Execution isomorphism invariance turns
these into observational equivalences. Compose the inverse left equivalence,
the contextual comparison, and the right equivalence. Since $N$ was arbitrary,
this proves the full existing all-allowances conclusion.

The input equivalence and the execution isomorphism law are explicit premises
of this inference theorem. They are not asserted for the proposition's programs
or its still-unconstructed copying execution.


\lean{Copying/TermContextCongruence}{26}{56}{TermContextCongruence.observationalEquiv_context_of_depth_eq}

\Needspace{8\baselineskip}
\subsection{Exact context depths and the remaining unequal-depth case}
\termdef{copying-context-depth}

Define $\delta_C(d)$ to be the context's outer allowance when its hole has
allowance $d$. A hole returns $d$, a lambda adds one, and a unary operation
retains the subcontext allowance. Application and binary contexts take the
maximum of the selected subcontext allowance and the fixed sibling's minimal
depth, in the original left/right order.

Induction on the context proves that this numerical function equals the
actual graph interpreter's depth: each constructor selects exactly the stated
successor, unchanged value or maximum. Applying that equality to a translated
term and then the exact constructor-interpretation identity gives
$\mathsf{depth}(C[t])=\delta_C(\mathsf{depth}(t))$.
These are statements about the unchanged interpreter, not a new translation.


\lean{Copying/ContextDepth}{10}{18}{ContextDepth.depthAt}

\lean{Copying/ContextDepth}{19}{25}{ContextDepth.depth_apply}

\lean{Copying/ContextDepth}{26}{32}{ContextDepth.depth_fill}

Three inequalities follow by context induction and elementary natural-number
arithmetic:
\[
 d\leq\delta_C(d),\qquad
 d\leq e\Longrightarrow\delta_C(d)\leq\delta_C(e),\qquad
 \delta_C(d+n)\leq\delta_C(d)+n.
\]
The hole cases are immediate. Lambda adds the same one to the inductive
comparison, and unary retains it. For a maximum-depth join the fixed sibling
is unchanged, so it cannot reverse monotonicity or increase the difference
beyond $n$. In particular, a fixed deep sibling may absorb an increase in the
hole allowance. These facts control the indices of the remaining graph-level
comparison, but do not themselves prove graph isomorphism.


\lean{Copying/ContextDepth}{33}{36}{ContextDepth.le_depthAt}

\lean{Copying/ContextDepth}{37}{42}{ContextDepth.depthAt_mono}

\lean{Copying/ContextDepth}{43}{49}{ContextDepth.depthAt_add_le}

For every ambient variable and atom scope, the private-cell term has minimal
depth three and the constant-result term has minimal depth two. Both equalities
hold by reducing the original syntax, its let/sequencing sugar and the original
translation depth function. Rewriting with these equalities proves that the
two depths differ. Thus the equal-depth theorem above cannot be applied
directly to this pair.

To remove the restriction, prove that interpreting $C$ with a graph $A$ at
$d$, or with $A$ padded to $e\geq d$, gives recursively isomorphic graphs
after both results are padded to a common allowance. This must retain the
original boxes, ordered argument exchange and distributors. The numeric
bounds establish adequate indices; the graph comparison remains unfinished.
Together with unrestricted filling and the cofinal criterion, it would remove
the equal-depth restriction. Appendix~B.1 retains this step, the concrete
execution laws, the copying UAM and the template proofs.

\lean{Copying/ContextDepth}{50}{52}{ContextDepth.privateCell_depth}

\lean{Copying/ContextDepth}{53}{56}{ContextDepth.constantResult_depth}

\lean{Copying/ContextDepth}{57}{62}{ContextDepth.proposition12_2_depths_differ}
